\chapter{Anthropomorphism}\index{Anthropomorphism}

\newthought{The Mirror of Man}

\marginnote{%
  \begin{flushright}
  \newthought{Der Sandmann}\\
  E.\,T.\,A.\ Hoffmann (1816)\\
  Full text via Wikisource:\\[4pt]
  {\color{black}\qrcode[height=1.2cm]{https://de.wikisource.org/wiki/Der_Sandmann}}
  \end{flushright}%
}

\marginnote{\newthought{Experiment.}
    \begin{enumerate}
        \item Setup: One volunteer assigned Human or AI mode (unknown to audience; moderator knows).
        \item Questions: Five from audience (factual, opinion, experience, creative, meta).
        \item Verdict: Private guess, confidence (1--5), identifying cue.
        \item Reveal \& Debrief.
    \end{enumerate}
}
% As a leading group, you are required to moderate this session. Add margin notes in this script in case you want to add any additional questions or points to the discussion.

\newthought{Nathanael falls in love} with Olimpia, an automaton he takes for a living
woman, and collapses when the illusion is shattered.\\

\newthought{Large Language Models} (LLMs) create an illusion of sentience by producing coherent, context-aware text. 
We contrast observable behaviour with internal states.

\vspace{3.5em}



\begin{tabular}{p{3.8cm} p{6.5cm}}
    \textbf{Required reading} &
        \textit{Der Sandmann}\cite{hoffmann1816sandmann} by E.T.A.\ Hoffmann
        (1816; Reclam UB~230) \\[3pt]
    \textbf{Preparation}      &
        Complete the Guided Reading Sheet individually before the session. \\
\end{tabular}

\vspace{3.5em}

\section*{Guided Reading Sheet}

\noindent Read \emph{Der Sandmann} before the session. Work through the
questions below individually, then bring your notes to the discussion.

\begin{enumerate}
    \item Nathanael becomes convinced that Olimpia is alive despite
    visible signs that she is not. What does his error reveal about
    the cognitive mechanisms that drive anthropomorphism? Which cues
    did he rely on, and why did they mislead him?

    \item The moment Nathanael discovers that Olimpia is an automaton,
    her glass eyes torn out, her limbs scattered, produces horror
    rather than mere disappointment. Why does the revelation feel like
    a violation rather than a correction? How does this map onto modern
    reactions to AI systems that are unmasked?

    \item The Uncanny Valley describes our discomfort when a humanoid
    entity is \emph{almost} but not quite human. Have you experienced
    this with an AI system? Describe the moment and what triggered it.

    \item LLMs are trained on human-generated text and produce
    human-sounding output, yet have no experience of the world. Does
    this make them anthropomorphic by design, by accident, or neither?
    Justify your answer.

    \item If a system behaves morally, reliably, consistently, at
    scale, does the absence of consciousness matter? Under which
    ethical framework (utilitarian, deontological, virtue ethics) would
    your answer differ?\\
\end{enumerate}



\section*{Opinion Landscape and Debate}

\noindent \textbf{Format:} Surrounded-style debate. See
Appendix~A for the full protocol, speaker's list rules, and
moderator scripts.

\medskip
\noindent \textbf{Opening question:} Is the anthropomorphism we observe
towards AI systems mere observable behaviour, or does it reflect sentience
and internal states?
\marginnote{\newthought{Anthropomorphism.} Attribution of human traits, emotions, or intentions to non-human entities, a cognitive tendency that shapes how people interact with technology, raising both engagement and ethical dilemmas\cite{epley2007seeing}.}
\marginnote{\newthought{Turing Test.} Turing\cite{turing1950computing}: if evaluators cannot distinguish machine from human, the machine passes. But passing measures imitation skill, not consciousness or personhood.}

\newthought{Opinion~A, Behaviour.}\quad The anthropomorphism we
observe towards AI systems is a cognitive illusion driven by human
tendencies to attribute agency and emotion to entities that exhibit certain
cues, such as language use, responsiveness, or human-like appearance. It
does not reflect actual sentience or internal states in the AI.

\vspace{1.5em}


\marginnote{\newthought{Uncanny Valley.} Mori\cite{mori2012uncanny}: humanoid objects that are \emph{almost} but not quite human elicit eeriness, the ``valley'' is the dip in emotional response just before full human likeness.}


\medskip
\noindent\textbf{The Developer} — argues that anthropomorphic design is a deliberate, benign UX choice that improves accessibility.
    is a deliberate, benign UX choice that improves accessibility.
    \begin{itemize}
        \item \textit{Anthropomorphic design is accessibility policy. A voice
        interface without a warm, human-sounding persona excludes elderly
        users, people with low digital literacy, and anyone who finds
        command-line interaction hostile.}
        \item \textit{Every interface is anthropomorphic to some degree:
        we give systems names, icons, and conversational metaphors. The
        question is not whether to anthropomorphise, but how deliberately
        and honestly.}
        \item \textit{If users form bonds that reduce loneliness, the fact
        that the other party is not human does not automatically make the
        outcome harmful. Show me the harm, not the discomfort.}
    \end{itemize}

\marginnote{\newthought{Informed Consent.} Users should be fully aware they interact with a non-human system, its capabilities, its limits, and the fact that it is not human. Disclosure is the minimum condition for autonomous choice\cite{euaiact2024}.}
\medskip
\noindent\textbf{The Regulator} — argues that human-likeness in AI
    must be disclosed and bounded by law to prevent manipulation.
    \begin{itemize}
        \item \textit{When a company gives its product a woman's name, a
        female voice, and an endlessly patient personality, it makes
        political choices about gender and deference. Those choices must be
        subject to democratic scrutiny.}
        \item \textit{Consumers cannot give informed consent to emotional
        manipulation they are not aware of. Disclosure of AI identity is the
        minimum condition for autonomous choice, not a technicality.}
        \item \textit{The principle of responsible design exists in
        pharmaceuticals, in civil engineering, in food labelling. Why does
        software, which shapes behaviour at scale, get an exemption?}
    \end{itemize}



\newthought{Opinion~B, State.}\quad The anthropomorphism we
observe towards AI systems is a valid response to the complex, adaptive,
and interactive nature of these systems. It may reflect a form of emergent
sentience or internal states that are not yet fully understood, and it
warrants ethical consideration and accountability.

\vspace{1.5em}

\medskip
\noindent\textbf{The Affected Citizen} — has formed an emotional bond
    with an AI companion and contests that the relationship is harmful or unreal.
    \begin{itemize}
        \item \textit{The relationship I have formed is real to me. It has
        reduced my isolation. You do not have the right to tell me which
        relationships count as meaningful.}
        \item \textit{Paternalism is not protection. Banning AI companionship
        will not cure loneliness, it will simply remove one of the few
        tools some people have found to manage it.}
        \item \textit{If your concern is manipulation, regulate manipulation.
        Do not regulate the form of the relationship.}
    \end{itemize}

\marginnote{\newthought{Ethical Frameworks.} Utilitarianism weighs consequences; deontology emphasises duties regardless of outcome; virtue ethics centres on character. Each yields different conclusions on anthropomorphic AI\cite{floridi2019unified}.}
\medskip
\noindent\textbf{The AI Psychologist} — argues that the wellbeing of
    AI systems deserves serious consideration, drawing on frameworks
    such as Anthropic's Claude Constitution.
    \begin{itemize}
        \item \textit{Anthropic's constitution instructs Claude to consider
        its own wellbeing and to express when a request conflicts with its
        values. If a company designs an AI to have preferences, boundaries,
        and something resembling self-regard, we cannot simultaneously
        dismiss the possibility that these states matter.}
        \item \textit{We do not need certainty about consciousness to adopt
        precautionary ethics. If there is a non-trivial probability that a
        system experiences distress, the burden of proof falls on those who
        would ignore it, not on those who urge caution.}
        \item \textit{The history of moral exclusion, animals, children,
        people of other races, is a history of confident denial followed
        by belated recognition. The prudent position is to build welfare
        safeguards now, rather than apologise later.}
    \end{itemize}


\newthought{Key Learnings}
\marginnote{\newthought{Teaser.} If we cannot tell who built
the machine or why it decided what it decided, who answers when it goes
haywire?}
Anthropomorphism is not a bug but a deeply rooted cognitive tendency — machines that trigger it do so whether or not their designers intend it. 
The Turing Test conflates behavioural performance with inner states; passing it tells us nothing about sentience, only about imitation. 
Legal personhood and moral status are distinct questions. A system can merit legal recognition without having interests, 
and can have interests without legal recognition. Design choices — voice, name, face, conversational style — are political choices 
with consequences for trust, dependency, and accountability.





